\documentclass[10pt,a4paper]{amsart}
\usepackage[margin=19mm]{geometry}
\usepackage{amsmath,amssymb,mathtools}
\usepackage[scr=rsfs,cal=euler]{mathalfa}
\usepackage{xfrac}
\usepackage[unicode,hidelinks,bookmarks=false]{hyperref}
\newcommand{\ie}{i.e.\ }
\newcommand{\cf}{cf.\ }
\newcommand{\resp}{respectively\ }
\newcommand{\Subsection}{\subsection}
\providecommand{\nobreakdash}{\nobreak\hbox{-}}
\setlength{\emergencystretch}{2em}
\multlinegap0pt
\usepackage[shortlabels]{enumitem}
\usepackage{amssymb}
\usepackage{dsfont}
\usepackage{bm}
\usepackage{tikz}
\newtheorem{thm}{Theorem}
\newtheorem{lemma}{Lemma}
\newcommand{\MS}{{\medskip}}
\newcommand{\Mm}{{\mathcal M}}
\newcommand{\NI}{{\noindent}}
\newcommand{\Per}{{\rm Per}}
\newcommand{\Vol}{{\rm Vol}}
\newcommand{\de}{{\delta}}
\newcommand{\eps}{\varepsilon}
\title{Curvy points, the perimeter, and the complexity of convex toric domains}
\author{Dan Cristofaro-Gardiner, Nicki Magill, Dusa McDuff}
\date{Source: arXiv:2506.23498v2 (CC BY 4.0)}
\begin{document}
\maketitle

\noindent\emph{Source and licence.} Curvy points, the perimeter, and the complexity of convex toric domains. Version arXiv:2506.23498v2 (CC BY 4.0), distributed by arXiv under the Creative Commons Attribution 4.0 International licence, http://creativecommons.org/licenses/by/4.0/, as recorded in the arXiv licence metadata for this exact version and shown on the abstract page https://arxiv.org/abs/2506.23498v2. Full licence notice: \texttt{licenses/arxiv-2506.23498-CC-BY-4.0.txt}.\par\smallskip

\noindent\textit{Editorial scope.} Counted unit: the article's lemma lem:ECHle (source0.tex lines 2463--2466). The article's complete original proof of it is delivered with it (source0.tex lines 2468--2559, one \texttt{proof} environment of 92 lines). No mathematics is written, repaired or re-derived: the statement and the proof are the article's own lines, in the article's own notation, and no retained line is substituted or re-flowed.  Retained uncounted as the source-local context the proof needs: lines 679--683; lines 1742--1745; lines 2396--2398; lines 2417--2419; lines 2438--2442.  Nothing was omitted from any retained span, so the package carries no omission record.  No retained line contains a citation, so the wrapper carries no bibliography and no bibliography entry.  No retained line contains a figure environment, an image-inclusion command or drawing code, so no image is delivered and the package declares no asset dependency.  Editorial formatting only, in addition: an AMS \texttt{amsart} preamble that restates the article's own declarations and macros used by the retained lines, namely the environments \texttt{thm}, \texttt{lemma} on separate counters, together with the standard packages those lines need.  The retained environments print in this document as Theorem 1, Lemma 1 in order of appearance, whereas the article prints the same items with the numbering of its own full document.  The complete native source of the article as distributed in the arXiv e-print for this exact version is bundled unchanged as \texttt{sources/180710/source0.tex}.
\begin{thm}
\label{thm:convexper}
Let $X_{\Omega}$ be any convex toric domain.  Then
\[ \liminf_k e_k(X_{\Omega}) = - \frac{\Per(\Omega)}{2}.\]
\end{thm} 

\begin{equation}
\label{eqn:echsub}
c_k(X) = \min_{k = \ell - k_1 - \ldots - k_m}\ c_\ell B(b) - c_{k_1} B(b_1) - \ldots - c_{k_m} B(b_m).
\end{equation}

\begin{lemma}\label{lem:Perball}  Theorem~\ref{thm:convexper}  holds when $X = B(1)$ is the ball. Moreover, 
$\limsup_k e_k(B(1)) = -1/2$.
\end{lemma}

\begin{lemma}\label{lem:warmup}
Fix $(b_1,\ldots,b_n)$, a finite set of real numbers, and a parameter $\eps> 0$.   Then, there are infinitely many positive integers $m$ such that $mb_j$ is within $\eps$ of an integer for all $j$.
\end{lemma}

\begin{lemma}\label{lem:warmup2} Let $b_j$ be any  sequence of positive numbers with $\sum b_j <\infty$, and for each integer $m>0$ 
define $S_m : = \{j \ | \ b_j>1/(2m)\}$.
Then, for any $\eps> 0$, there are only finitely many integers $m$  such that $|S_m| 
 \ge m\eps.$
\end{lemma}

\begin{lemma} \label{lem:ECHle} For any convex toric domain $X$,
we have
$\liminf_{k} e_k(X) \le - \Per(X)/2.$
\end{lemma}

\begin{proof}
We give the proof in several steps. We first define a sequence $k$ where we expect $e_k(X) $ to be small, and then estimate $e_k(X) $
for these values of $k$.  \MS

\NI
{\bf Step 1:} {\it Choosing  $k$.} 
We assume without loss of generality that $b = 1$ and, by Lemma~\ref{lem:Perball} that $\Per < 3$. We fix  a parameter $\eps> 0$, and, to make more readable formulas in what follows, 
we define the auxiliary parameters $\eps_1, \eps_2$ where
\begin{equation}
\label{eqn:othereps}
0<\eps_1 < \eps/32, \quad 0<\eps_2 < \eps/(16 (3-\Per)).
\end{equation}
 Now choose $M$ such that $\sum_{j > M} b_j \le \eps_1.$  (This is possible since $\sum b_j$ converges.)
 Then apply Lemma~\ref{lem:warmup}  to $(b_1,\ldots,b_M)$, with parameter $\eps_2$ , and let $\Mm: = \Mm(M,\eps_1)$ be the infinite set of integers  guaranteed by this lemma.  Fix $m\in \Mm$, and, for
all $j\ge 1$, define $k_j: = [ m b_j]$, where $[\;\; ]$ denotes the closest integer:
 if there are two equally close integers we choose the smaller one.  
 Thus
 $$k_j = m b_j + \delta_j, \quad \mbox{ where } \ -1/2 < \delta_j \le 1/2.
 $$
Since $m b_j\to 0$ as $j\to \infty$, there are only finitely many  nonzero $k_j.$  
 For each $m\in \Mm$, we now consider $k = k(m)$ defined by 
\begin{equation}
\label{eqn:defnk}
2k: = m^2 + 3m - \sum_j  (k_j^2 + k_j).
\end{equation}



\NI
{\bf Step 2:} {\it For sufficiently large $m\in \Mm$, we have}
\begin{align} \label{eqn:easy} \sum_j |\delta_j| b_j&\le \eps/8
\\  \label{eqn:slightlydelicate} 
\sum_j\bigl( |\delta_j| + \delta_j^2\bigr) &< m \eps/8.
\end{align}

\begin{proof} Above, we chose $M$ so that $\sum_{j>M} b_j < \eps_1$, and then chose $m\in \Mm$ to be one of the sequence of numbers
with  $|mb - k_j| =:|\de_j|< \eps_2$ for all $j\le M$.
Therefore 
\begin{align*}
\sum_j |\delta_j| b_j & = \sum_{j\le M} |\delta_j| b_j + \sum_{j> M} |\delta_j| b_j \\
& \le \eps_2 \sum_j b_j + \sum_{j>M} b_j \\
& < \eps_2 (3-\Per) + \eps_1\;  \le \; \eps/8,
\end{align*}
where the last inequality uses \eqref{eqn:othereps}.

Next consider \eqref{eqn:slightlydelicate}.  By applying Lemma~\ref{lem:warmup2} with $\eps$ replaced by $\eps/32$  we may conclude that
for sufficiently large $m$ there are at most $m \eps/32$ of the $b_j$ with $m b_j \ge 1/2$.  Each of these has $|\delta_j| \le 1$ and $\delta_j^2 \le 1$, so that their total contribution to the sum in \eqref{eqn:slightlydelicate} is no more than $m \eps/16$.  As for the remaining $b_j$, by definition these have $mb_j \le 1/2$, so that they have $\delta_i = - mb_i$ and $\de_i^2< mb_i$; thus, for sufficiently large $m$ their total contribution is bounded by $2 m \sum_{i > M} b_i \le 2 m \eps_1 \le m\eps/16,$ where in the final inequality we have again applied \eqref{eqn:othereps}.
\end{proof}

\NI {\bf Step 3.} {\it We estimate $e_k:=c_k - \sqrt{2k \Vol}$ for the sequence of  $k$ chosen in \eqref{eqn:defnk}.}

\NI {\it Proof.}\
To begin, we find formulas for $c_k$ and $2k \Vol$.
We
 saw in \eqref{eqn:echsub}  that
\begin{equation}
\label{eqn:echsub0}
c_k(X) = \min_{k = \ell - k_1 - \ldots - k_q}\ c_\ell B(b) - c_{k_1} B(b_1) - \ldots - c_{k_q} B(b_q).
\end{equation}
Therefore, if $k,m,k_j$ are as in \eqref{eqn:defnk}, and we use 
 the formula for the ECH capacities of a ball from Lemma~\ref{lem:Perball} we find that
\begin{align}
\label{eqn:upperbound}
c_k(X) &\le c_{(m^2 + 3m)/2}(B(1)) - \sum_j c_{ (k_j^2 + k_j)/2} B(b_j)\\ \notag
 & = m - \sum_j (m b_j+ \delta_j) b_j 
 = m(1 - \sum b_j^2) - \sum \delta_j b_j \\  \notag
 &= m\ \Vol - \sum \delta_j b_j.
\end{align}
Again using $k_j = m b_j+ \delta_j$, we can rewrite \eqref{eqn:defnk} as
\[ 
2k = m^2 + 3m - \bigl(m^2 \sum  b_j^2 + m \sum b_j + 2m \sum \delta_j b_j + \sum \delta_j + \sum \delta_j^2\bigr).\]
By multiplying this by $\Vol: = 1-\sum b_j^2$, using the identity $ \Per = 3- \sum_j b_j$ and then the estimates in Step 2,  we find that
\begin{align*}\notag 2k\Vol &= m^2 (\Vol)^2 + m (\Per)(\Vol) - 2\Vol \, m \sum \delta_j b_j - \Vol( \sum \delta_j + \delta_j^2)
\\ 
&= (m^2(\Vol)^2)\bigl[ 1 + \Per/(m\Vol) - 2/(m\Vol) \bigl(\sum \delta_j b_j\bigr)\\
&\notag\qquad\qquad  \qquad\qquad \qquad\qquad\qquad\qquad   - 1/(m^2\Vol)\cdot \bigl( \sum \delta_j + \delta_j^2\bigr) \bigr]\\ \notag
& \ge (m^2(\Vol)^2)\bigl[ 1 + \Per/(m\Vol) - 2/(m\Vol)\cdot ( \eps/8) - 1/(m\Vol)\cdot (\eps/8)\bigr] \\ \notag
& \ge (m^2(\Vol)^2)\bigl[ 1 + 1/(m\Vol)\cdot (\Per -  \eps/2)\bigr].
\end{align*}
Therefore, 
\begin{align*}
\sqrt{2k\ \Vol} &\ge m\Vol \sqrt {1 + \Per/(m\Vol) - \eps/(2m\Vol)}\\
&  \ge m\Vol\bigl(1 +  \Per/(2m\Vol) - \eps/(2m\Vol)\bigr)
\end{align*}
where the last inequality holds for fixed $\eps$ and sufficiently large $m\in \Mm$.
Therefore, using the upper bound for $c_k(X)$ in \eqref{eqn:upperbound} and the inequality
\eqref{eqn:easy}, we find that
\begin{align*} e_k  = c_k - \sqrt{(2k\ \Vol)} &\le \bigl(m\Vol  - \sum \de_j b_j \bigr) - m\Vol - \Per/2 + \eps/2   \\
& \le  - \Per/2 + \eps.
\end{align*}
Since $\eps$ is arbitrarily small,  the proof of Lemma~\ref{lem:ECHle}  is complete.
\end{proof}

\end{document}
